\documentclass[11pt]{article}
\usepackage[margin=1in]{geometry}
\usepackage[utf8]{inputenc}
\usepackage{amsmath,amssymb}
\usepackage{graphicx}
\usepackage{booktabs}
\usepackage{caption}
% Left-justify every caption; without singlelinecheck=false LaTeX
% centres any caption short enough to fit one line, so short and long
% captions in the same paper are set differently.
\captionsetup{justification=raggedright,singlelinecheck=false}
\usepackage[hidelinks]{hyperref}
\usepackage{url}
\usepackage{enumitem}
\usepackage[round,sort,authoryear]{natbib}

\newcommand{\dkl}{$\Delta$KL}

\title{\bfseries Vocabulary position and trademark lifecycles:\\
\large An event-dated corpus and a lead/lag text measure for the commercial economy}
\author{Eric Silver\thanks{Independent researcher; formerly a PhD candidate at Carnegie Mellon University. Comments to \texttt{epsilver@gmail.com}. The author's current employment is unrelated to this work, and the views expressed are the author's alone. I thank Simon DeDeo for the topic-scoring suggestion recorded in Section~\ref{sec:construct}, and Carolina Castaldi for comments on an earlier draft and for orienting me in the trademark-indicator literature. Neither has seen or endorsed the present version, and all errors and judgements are mine.}}
\date{July 2026 working paper}

\begin{document}
\maketitle

\begin{abstract}
\noindent All 13.99
million USPTO trademark case files are re-parsed here into event-dated
prosecution histories, and every filing scored on how surprising its language
is against what its industry filed in the five years before it and in the five
years after. The average of the two surprises is the filing's
\emph{atypicality}; their difference is its \emph{lead}, positive for wording
the industry subsequently moved toward (\emph{leading}) and negative for
wording it was already leaving behind (\emph{lagging}).

\noindent Within the trademark record, \emph{leading} marks are less likely to
remain in commerce. Of the most leading fifth of 3.10 million registrations,
46.7\% survive the sworn re-proof of use that a US registration requires
between its fifth and sixth year, against 48.1\% of the most lagging fifth ---
a gap of 1.4 percentage points, within Nice class and registration year, with owner-clustered
errors and drafting and filer controls ($t = 8.3$). The gap holds across
industries and across
filer types, but not across time: it was two to three times larger for marks
filed in the late 1990s, reversed for those filed 2000--2004, and back from
2008, with the entire swing inside four technology classes and a flat one to
four points everywhere else. Within those classes, most of the swing is which
themes each era's leading filings were drawn from; inside a theme, the
penalty is a steadier two to five points.

\noindent 0.83\% of first-time filers with professional representation ever
appear in SEC financial-reporting records, against 0.18\% of self-filers. The
atypicality of a firm's first registered mark does not predict whether its
owner ever reaches SEC reporting once class, year and description length are
absorbed; yet both leading ($\mu = 0.39$\%) and lagging ($\mu = 0.44$\%)
owners outperform the persistent middle ($\mu = 0.24$\%). Neither axis is a proxy for
patenting: across 134{,}075 matched firm-years, lead correlates with
$\log(1+\mathrm{patents})$ at $r = -0.020$ and atypicality at $r = +0.009$,
and within a firm over time at $-0.043$ and $+0.011$.

\noindent Atypicality can be measured on terms or on themes, and the two
disagree. Term atypicality predicts both outcomes negatively ($-0.09$pp of SEC
reporting per standard deviation of the score; $-0.055$ standard deviations
of log patenting); theme atypicality does not predict SEC reporting and
predicts patenting weakly the other way ($+0.012$ standard deviations). The two therefore disagree in sign on both outcomes, which
makes a text measure's sign at the firm level a property of how the text was
counted rather than of the firm. Every estimate is conditional on the office
having granted the mark, because registration is itself selected on language;
without that conditioning the same data support weaker and sometimes opposite
conclusions. Code, panels, and the corpus build are public at
\url{https://github.com/ericsilver/tm-vocabulary}, with a browsable explorer
for the per-theme and per-Nice-class results behind the pooled estimates.
\end{abstract}

\section{Introduction}

Text-based measurement has reshaped how economists quantify invention.
\citet{Kelly2021} score every US patent since 1840 by its textual
distance from prior and similarity to subsequent patents, and
\citet{HobergPhillips2016} build product-market industries from 10-K text;
\citet{Kalyani2024} measures patent creativity as the share of technical
terminology not seen before and ties it to firm productivity. Both
programs inherit the patent record's coverage: technical invention by
the small minority of firms that patent. Most firms never patent,
service-sector and business-model innovation rarely patents, and
patenting propensity varies so much across industries that
cross-industry comparisons are fraught \citep{Graham2013, Dinlersoz2018}.
Trademark filings offer the complementary window. Nearly every
commercially serious product or service introduction in the United
States is accompanied by a trademark filing, the filing's goods/services
description states in plain commercial language what the firm intends to
sell, and the corpus extends across every sector of the economy at the
rate of hundreds of thousands of granted filings per year.

The trademark--innovation literature is three decades old and has moved
well past counting \citep{Castaldi2020}. \citet{Flikkema2025} survey its first ten years of
validation work and its open questions; \citet{Willeke2026} inventory the
data sources and the preprocessing burden that has kept them underused.
Micro-level links to firm dynamics are established \citep{Dinlersoz2018,
Dinlersoz2023}, and the goods/services description is already read as
text rather than treated as a label: \citet{Semadeni2006} scores
consulting service-mark descriptions into a competitive positioning
space, \citet{SemadeniAnderson2010} tracks whether one firm's described
offering is imitated by rivals, and \citet{Flikkema2019} codes
descriptions against the innovation each mark accompanies. Prosecution
records have been exploited too: \citet{Fink2024} use contested filing
sequences to study trademark races, and their outcome --- continued use
five years after registration --- is the margin this paper's central
result also runs on. What \citet{Flikkema2025} put as the field's next
step is the move from counts and innovation intensity to the \emph{type},
\emph{quality}, and \emph{trend} content of what marks cover: the
trademark analogue of what \citet{Kelly2021} and \citet{Kalyani2024} did
for patent text. That analogue does not yet exist at corpus scale, and
the two strands above --- description text and prosecution events --- have
run on separate data.

This paper builds both on one corpus, so that what a filing says and what
subsequently happened to it are observed on the same record.

Whether a description's language is strange for its industry predicts very
little; whether it arrived \emph{before} its industry's language moved that
way predicts a good deal. Marks whose wording the industry subsequently
adopted are abandoned at the five-year use-proof gate more often than marks
written in language the industry was already leaving behind. Their owners
also reach SEC financial reporting less often than the owners of the most
lagging marks, though both tails do better than the middle.

The USPTO distributes its trademark backfile
as bulk XML --- 83 archives covering 1884 to 2025 --- from which the
published research files retain a filing's current status but not the dated
sequence of events that produced it. Re-parsing the backfile yields 242
million prosecution events --- every step in the examination of an
application --- across 13.99 million case files, each event
carrying a code and a date: every office action, response, maintenance
filing, and cancellation. The event codes are resolved into a 760-code dictionary. A header pass
extracts counsel of record, the statutory filing basis, and the
postregistration declarations. A third pass recovers owner domicile, which
the class-level files omit. Owner names are then normalized to a clustering
key and resolved to two external universes: SEC registrants, through the
agency's own company-name and ticker files, and Regulation~D issuers,
through the Form~D quarterly data sets --- the latter filtered to operating
companies, with ambiguous names dropped rather than guessed. That link is
what lets a mark's language be followed to financing and to public listing.
Where the patent record dates invention once, at grant, the trademark
prosecution record dates a commercial claim's entire life, including the
use-proof gates at which marks attached to products that never materialized
are culled.

Every filing is grouped by its industry class --- one of the 45
international goods and services categories set by the Nice Agreement ---
and placed by its filing
date, then put to two readers who have seen disjoint halves of that
industry's stream: one who knows only what the class filed before it, one
who knows only what it filed after. Each reader's ``how surprising'' is a
Kullback--Leibler divergence between the filing's mix of topics --- its
weights over word-clusters fitted to the corpus --- and the industry's,
large when the filing puts weight where its industry's filings put little.
\emph{Atypicality} is the average of the two answers. \emph{Lead}, written
\dkl{}, is the gap. Language that surprised the first reader and not the
second was surprising because it came first: the industry moved toward it.
Language that read as ordinary to the first and unusual to the second was
left behind as the field moved on. Filings of the first kind are
\emph{leading}, of the second \emph{lagging}.

The construct adapts \citet{Murdock2017} and \citet{Barron2018};
Section~\ref{sec:construct} sets out what transfers and what does not. The
nearest measurement precedent on this corpus is \citet{vonGraevenitz2022},
who flag goods/services tokens that are both new and fast-spreading; their
unit is the token and their flag binary, while the score here is continuous
and attaches to the filing, so every application carries a position that can
be crossed with its own prosecution outcomes.

Results are stated on marks the office granted. Registration is itself
selected on language, so a correlation over all filings mixes what
examination does to a description with what the market does to the product;
Appendix~\ref{app:uncond} reports the unconditional version. Two further
products of the exercise are released with it: the corpus, and a
measurement hazard that applies to text-based novelty measures generally.

